% from ega4-14-fr.tex
\oldpage[IV-3]{214}

  \item[a$'$)] La propriété $P(Y')$ est vraie pour tout morphisme fini
  $g:Y'\to Y$.
  \item[a$''$)] La propriété $P(Y'')$ est vraie pour le normalisé $Y''$ de
  $Y_{\mathrm{red}}$ (\textbf{II}, 6.3.8).
  \item Pour tout point $x''$ de $X''=X\times_Y Y''$ au-dessus de $x$, il existe
  une composante irréductible $Z''$ de $X''$ contenant $x''$ et
  équidimensionnelle sur $Y''$ au point $x''$.
\end{enumerate}
\end{proposition}

Il est trivial que $a)$ implique $a')$. Pour montrer que $a')$ implique $a'')$
notons que l'on peut écrire $Y''=\operatorname{Spec}(\mathcal B)$, où $\mathcal B$
est une $\mathcal O_Y$-Algèbre quasi-cohérente entière sur $\mathcal O_Y$~; comme
$Y$ est noethérien, $\mathcal B$ est limite inductive de ses sous-$\mathcal O_Y$-
Algèbres $\mathcal B_\lambda$ qui sont quasi-cohérentes et de type fini (I, 9.6.6)~;
mais alors les $\mathcal B_\lambda$ sont des $\mathcal O_Y$-Algèbres \emph{finies}
(\textbf{II}, 6.1.2)~; on peut donc écrire
$Y''=\varprojlim Y'_\lambda$, où
$Y'_\lambda=\operatorname{Spec}(\mathcal B_\lambda)$, d'où
$X''=X\times_Y Y''=\varprojlim X'_\lambda$, avec
$X'_\lambda=X\times_Y Y'_\lambda$. En vertu de $a')$, les morphismes
$f'_\lambda:X'_\lambda\to Y'_\lambda$ sont ouverts en tous les points de
$X'_\lambda$ au-dessus de $x$~; on en conclut que $f''$ est ouvert en tous les
points de $X''$ au-dessus de $x$, par (8.10.1) et (14.3.3.1, (i)).

Comme le préschéma $Y''$ est \emph{normal} par définition, le fait que $b)$ entraîne
$a'')$ résulte de (14.4.4) appliqué à la composante irréductible
équidimensionnelle de l'énoncé et à la restriction de $f''$ à cette composante. Il
reste donc à montrer que $a'')$ entraîne $a)$ et $b)$. Compte tenu de (1.10.3), on
peut se borner au cas où $Y=\operatorname{Spec}(\mathcal O_y)$, en notant que le
morphisme canonique $\operatorname{Spec}(\mathcal O_y)\to Y$ est universellement
bicontinu (I, 3.6.5), et d'autre part que
$Y''\times_Y\operatorname{Spec}(\mathcal O_y)$ est le normalisé de
$(\operatorname{Spec}(\mathcal O_y))_{\mathrm{red}}$ comme il résulte de la
permutabilité des opérations de fermeture intégrale et de localisation (Bourbaki,
\emph{Alg. comm.}, chap.~V, \S~1, n\textsuperscript{o}~5, prop.~16). Supposant donc
$Y=\operatorname{Spec}(A)$, où $A$ est un anneau local noethérien, et $y$ le point
fermé de $Y$, on sait ($\mathbf{0}_{\mathrm{IV}}$, 23.2.5) qu'il existe une
factorisation
\[
  Y''\xrightarrow{u}Y_1\xrightarrow{v}Y
\]
du morphisme structural, telle que $v$ soit un morphisme \emph{fini surjectif}, $u$
un morphisme entier, \emph{radiciel et dominant}, donc surjectif (\textbf{II},
6.1.10) et par suite un \emph{homéomorphisme universel} (2.4.5). Si l'on pose
$X_1=X\times_Y Y_1$, la projection $X''\to X_1$ est donc un homéomorphisme, et
l'hypothèse $a'')$ entraîne par suite que $f_1=f_{(Y_1)}:X_1\to Y_1$ est ouvert en
tous les points de $X_1$ au-dessus de $x$. En outre, cela montre que pour prouver
la propriété $b)$, il suffit de prouver la même propriété où l'on remplace $Y''$,
$X''$ et $x''$ par $Y_1$, $X_1$ et un point $x_1$ de $X_1$ au-dessus de $x$~: Mais
$Y_1$ est noethérien et en outre il est géométriquement unibranche puisque $Y''$
est normal et $u$ radiciel (6.15.1)~; la propriété à prouver résulte donc de
(14.4.1). Il reste à montrer que $f$ est universellement ouvert au point $x$, ce
qui résultera du lemme suivant~:

\begin{lemma}[14.4.8.1]
\label{IV.14.4.8.1-fr}
Soit $v:Y_1\to Y$ un morphisme fermé (resp. universellement fermé) et surjectif.
Pour qu'un morphisme $f:X\to Y$ soit ouvert (resp. universellement ouvert) en un
point $x\in X$, il suffit que $f_1=f_{(Y_1)}:X_1=X\times_Y Y_1\to Y_1$ soit ouvert
(resp. universellement ouvert) en tous les points de $X_1$ au-dessus de $x$.
\end{lemma}

La seconde assertion résulte trivialement de la première et du fait que pour tout
changement de base $Y'\to Y$, le morphisme
$v_{(Y')}:Y_1\times_Y Y'\to Y'$ est encore surjectif et est fermé si $v$ est
universellement fermé. Pour prouver la première assertion, considérons

